% Part: incompleteness
% Chapter: arithmetization-syntax
% Section: proofs-in-nd

\documentclass[../../../include/open-logic-section]{subfiles}

\begin{document}

\olfileid{inc}{art}{pnd}
\olsection{సహజ నిగమనంలో \usetoken{P}{derivation}}

\begin{explain}
\tetoken{వ్యుత్పత్తులను}{!!{derivation}s} అంకగణితీకరించాలంటే,
\tetoken{వ్యుత్పత్తులను}{!!{derivation}s} సంఖ్యలుగా ప్రాతినిధ్యం చేయాలి.
\tetoken{వ్యుత్పత్తులు}{!!{derivation}s}
\tetoken{సూత్రాల}{!!{formula}s} వృక్షాలు, ప్రతి నిగమనానికీ ఒకటి లేదా రెండు
గుర్తింపు చీటీలు ఉంటాయి కనుక, పునరావృత్త ప్రాతినిధ్యమే అత్యంత సహజమైన
పద్ధతి. ఒక \tetoken{వ్యుత్పత్తిని}{!!{derivation}} చివరి నిగమనపు
పూర్వాధారాలకు దారితీసే తక్షణ ఉప-\tetoken{వ్యుత్పత్తుల}{!!{derivation}s}
సంఖ్య, ఆ ఉప-\tetoken{వ్యుత్పత్తుల}{!!{derivation}s} ప్రాతినిధ్యాలు,
అంత్య-\tetoken{సూత్రం}{!!{formula}}, చివరి నిగమనపు ఉపసంహరణ చీటీ, చివరి
నిగమనం రకాన్ని సూచించే సంఖ్య అనే భాగాలు గల క్రమిత సమూహంగా ప్రాతినిధ్యం
చేస్తాం.
\end{explain}

\begin{defn}
  $\delta$ సహజ నిగమనంలో \tetoken{ఒక వ్యుత్పత్తి}{!!a{derivation}} అయితే,
  $\Gn{\delta}$ను ఆగమనాత్మకంగా ఇలా నిర్వచిస్తాం:
  \begin{enumerate}
  \item
    $\delta$లో పరికల్పన~$!A$ మాత్రమే ఉంటే, $\Gn{\delta}$ అనేది
    $\tuple{0,\Gn{!A},n}$. అది \tetoken{ఉపసంహరించని}{!!{undischarged}}
    పరికల్పన అయితే సంఖ్య~$n$ అనేది~$0$; లేకపోతే అది సంఖ్యా చీటీ.
  \item
    $\delta$ సున్నా, ఒకటి, రెండు లేదా మూడు పూర్వాధారాలు గల నిగమనంతో ముగిస్తే,
    వరుసగా $\Gn{\delta}$ అనేది
    \begin{align*}
      & \tuple{0,\Gn{!A},n,k},\\
      & \tuple{1,\Gn{\delta_1},\Gn{!A},n,k},\\
      & \tuple{2,\Gn{\delta_1},\Gn{\delta_2},\Gn{!A},n,k},\text{ లేదా}\\
      & \tuple{3,\Gn{\delta_1},\Gn{\delta_2},\Gn{\delta_3},\Gn{!A},n,k},
    \end{align*}
    ఇక్కడ $\delta_1$, $\delta_2$, $\delta_3$లు $\delta$లోని చివరి నిగమనపు
    పూర్వాధారం లేదా పూర్వాధారాల వద్ద ముగిసే
    ఉప-\tetoken{వ్యుత్పత్తులు}{!!{derivation}s}; $!A$ అనేది $\delta$లోని ఆ చివరి నిగమనపు
    నిష్కర్ష; $n$ చివరి నిగమనపు ఉపసంహరణ చీటీ (నిగమనం ఏ పరికల్పననూ
    ఉపసంహరించకపోతే $0$); చివరి నిగమనంలో వాడిన నియమాన్ని బట్టి $k$ను ఈ
    పట్టిక ఇస్తుంది.

    \begin{tabular}{lccccccc}
      \text{నియమం:} & \Intro{\land} & \Elim{\land} & \Intro{\lor} & \Elim{\lor} \\
      $k$: & 1 & 2 & 3 & 4 \\[1.5ex]
      \text{నియమం:} & \Intro{\lif} & \Elim{\lif} & \Intro{\lnot} & \Elim{\lnot} \\
      $k$: & 5 & 6 & 7 & 8 \\[1.5ex]
      \text{నియమం:} & \FalseInt & \FalseCl & \Intro{\lforall} & \Elim{\lforall} \\
      $k$: & 9 & 10 & 11 & 12 \\[1.5ex]
      \text{నియమం:} & \Intro{\lexists} & \Elim{\lexists} & \Intro{\eq} & \Elim{\eq} \\
      $k$: & 13 & 14 & 15 & 16
    \end{tabular}
  \end{enumerate}
\end{defn}

\begin{ex}
  ఈ అతి సరళమైన \tetoken{వ్యుత్పత్తిని}{!!{derivation}} పరిగణించండి:
  \begin{prooftree}
    \AxiomC{$\Discharge{!A \land !B}{1}$}
    \RightLabel{\Elim{\land}}
    \UnaryInfC{$!A$}
    \DischargeRule{\Intro{\lif}}{1}
    \UnaryInfC{$(!A \land !B) \lif !A$}
  \end{prooftree}
  పరికల్పన గ్యోడెల్ సంఖ్య $d_0=\tuple{0,\Gn{!A\land !B},1}$.
  $\Elim{\land}$ నిష్కర్షతో ముగిసే
  \tetoken{వ్యుత్పత్తి}{!!{derivation}} గ్యోడెల్ సంఖ్య
  $d_1=\tuple{1,d_0,\Gn{!A},0,2}$ అవుతుంది
  ($\Elim{\land}$కు ఒక పూర్వాధారం ఉంది కనుక $1$; నిష్కర్ష~$!A$ యొక్క
  గ్యోడెల్ సంఖ్య; ఏ పరికల్పనా ఉపసంహరించబడలేదు కనుక $0$;
  $\Elim{\land}$ను సంకేతీకరించే సంఖ్య $2$). అప్పుడు మొత్తం
  \tetoken{వ్యుత్పత్తి}{!!{derivation}} గ్యోడెల్ సంఖ్య
  $\tuple{1,d_1,\Gn{((!A\land !B)\lif !A)},1,5}$, అంటే
  \[
  \tuple{1,\tuple{1,\tuple{0,\Gn{(!A\land !B)},1},\Gn{!A},0,2},
    \Gn{((!A\land !B)\lif !A)},1,5}.
  \]
\end{ex}

\begin{explain}
\tetoken{వ్యుత్పత్తుల}{!!{derivation}s} ప్రాతినిధ్యాన్ని నిర్ణయించిన తరువాత,
అటువంటి \tetoken{వ్యుత్పత్తుల}{!!{derivation}s} గ్యోడెల్ సంఖ్యలను ఆదిమ
పునరావృత్తంగా మార్చగలమనీ, వాటి ముఖ్య ధర్మాలు, సంబంధాలను అలాగే వ్యక్తీకరించగలమనీ
చూపాలి. కొన్ని క్రియలు సరళమైనవి: ఉదాహరణకు,
\tetoken{ఒక వ్యుత్పత్తి}{!!a{derivation}} గ్యోడెల్ సంఖ్య~$d$ ఇస్తే,
$\fn{EndFmla}(d)=(d)_{(d)_0+1}$ దాని అంత్య-\tetoken{సూత్రం}{!!{formula}}
గ్యోడెల్ సంఖ్యను, $\fn{DischargeLabel}(d)=(d)_{(d)_0+2}$ ఉపసంహరణ
చీటీని, $\fn{LastRule}(d)=(d)_{(d)_0+3}$ చివరి నిగమనం రకాన్ని సూచించే
సంఖ్యను ఇస్తాయి. కొన్ని క్రియలు మరింత కష్టం. వాటిని ఎలా చేయాలో కనీసం
రూపురేఖలుగా చూద్దాం. ``$\delta$ అనేది $\Gamma$ నుంచి $!A$ యొక్క
\tetoken{ఒక వ్యుత్పత్తి}{!!a{derivation}}'' అనే సంబంధం $\delta$, $!A$ల
గ్యోడెల్ సంఖ్యల ఆదిమ పునరావృత్త సంబంధమని చూపడమే లక్ష్యం.
\end{explain}

\begin{prop}
  క్రింది సంబంధాలు ఆదిమ పునరావృత్తం:
  \begin{enumerate}
  \item $!A$ అనేది $\delta$లో చీటీ~$n$తో పరికల్పనగా కనిపిస్తుంది.
  \item $\delta$లో చీటీ~$n$ గల పరికల్పనలన్నీ $!A$ రూపంలోనే ఉన్నాయి
    (అంటే $\delta$లో చీటీ~$n$తో పరికల్పన~$!A$ను
    \tetoken{ఉపసంహరించవచ్చు}{!!{discharge}}).
  \end{enumerate}
\end{prop}

\begin{proof}
  \tetoken{సూత్రాల}{!!{formula}s} గ్యోడెల్ సంఖ్యలకూ,
  \tetoken{వ్యుత్పత్తుల}{!!{derivation}s} గ్యోడెల్ సంఖ్యలకూ మధ్య అనురూపమైన
  సంబంధాలు ఆదిమ పునరావృత్తమని చూపాలి.
  \begin{enumerate}
  \item గ్యోడెల్ సంఖ్య~$d$ గల \tetoken{వ్యుత్పత్తిలో}{!!{derivation}}
    చీటీ~$n$ గల పరికల్పన గ్యోడెల్ సంఖ్య $x$ అయినప్పుడు నిలిచే
    $\fn{Assum}(x,d,n)$ ఆదిమ పునరావృత్తమని చూపాలి. గ్యోడెల్ సంఖ్య
    $\tuple{0,x,n}$ గల \tetoken{వ్యుత్పత్తి}{!!{derivation}} అనేది $d$కు
    ఉప-\tetoken{వ్యుత్పత్తి}{!!{derivation}} అయినప్పుడే ఇది జరుగుతుంది.
    మన \tetoken{వ్యుత్పత్తుల}{!!{derivation}s} సంకేతీకరణ
    \olref[cmp][rec][tre]{sec}లో ప్రవేశపెట్టిన వృక్ష సంకేతీకరణకు ఒక ప్రత్యేక
    సందర్భం. కాబట్టి ఆదిమ పునరావృత్త ప్రమేయం $\fn{SubtreeSeq}(d)$,
    $d$ యొక్క అన్ని ఉప-\tetoken{వ్యుత్పత్తుల}{!!{derivation}s} గ్యోడెల్
    సంఖ్యల క్రమాన్ని ఇస్తుంది (దాని పొడవు గరిష్ఠంగా $d$). అందువల్ల
    \[
    \fn{Assum}(x,d,n)\defiff\bexists{i<d}{(\fn{SubtreeSeq}(d))_i=
      \tuple{0,x,n}}.
    \]
    అని నిర్వచించవచ్చు.
  \item గ్యోడెల్ సంఖ్య~$d$ గల \tetoken{వ్యుత్పత్తిలో}{!!{derivation}}
    చీటీ~$n$ గల పరికల్పనలన్నీ గ్యోడెల్ సంఖ్య~$x$ గల
    \tetoken{సూత్రమే}{!!{formula}} అయితే నిలిచే $\fn{Discharge}(x,d,n)$
    ఆదిమ పునరావృత్తమని చూపాలి. ఈ సంబంధం
    $\bforall{y<d}{(\fn{Assum}(y,d,n)\lif y=x)}$ అయినప్పుడే నిలుస్తుంది.
  \end{enumerate}
\end{proof}

\begin{prop}
  \ollabel{prop:followsby}
  గ్యోడెల్ సంఖ్య~$d$ గల \tetoken{వ్యుత్పత్తి}{!!{derivation}}~$\delta$లోని
  చివరి నిగమనం సరైనదైతేనే నిలిచే $\fn{Correct}(d)$ ధర్మం ఆదిమ పునరావృత్తం.
\end{prop}

\begin{proof}
  ప్రతి నిగమన నియమం~$R$కి $\fn{FollowsBy}_R(d)$ సంబంధం ఆదిమ పునరావృత్తమని
  చూపాలి. $d$ అనేది \tetoken{వ్యుత్పత్తి}{!!{derivation}}~$\delta$ గ్యోడెల్
  సంఖ్యగా ఉండి, $\delta$ అంత్య-\tetoken{సూత్రం}{!!{formula}} దాని తక్షణ
  ఉప-\tetoken{వ్యుత్పత్తుల}{!!{derivation}s} నుంచి $R$ను సరిగా
  వర్తింపజేయడం వల్ల వచ్చినప్పుడే $\fn{FollowsBy}_R(d)$ నిలుస్తుంది.

  సరళమైన సందర్భం \Intro{\land} నియమం. $\delta$ సరైన
  $\Intro{\land}$ నిగమనంతో ముగిస్తే అది ఇలా ఉంటుంది:
  \begin{prooftree}
    \AxiomC{}
    \RightLabel{$\delta_1$}
    \DeduceC{$!A$}

    \AxiomC{}
    \RightLabel{$\delta_2$}
    \DeduceC{$!B$}

    \RightLabel{\Intro\land}
    \BinaryInfC{$!A\land !B$}
  \end{prooftree}
  అప్పుడు $\delta$ గ్యోడెల్ సంఖ్య~$d$ అనేది
  $\tuple{2,d_1,d_2,\Gn{(!A\land !B)},0,k}$; ఇక్కడ
  $\fn{EndFmla}(d_1)=\Gn{!A}$, $\fn{EndFmla}(d_2)=\Gn{!B}$, $n=0$,
  $k=1$. కాబట్టి $\fn{FollowsBy}_{\Intro\land}(d)$ను ఇలా నిర్వచించవచ్చు:
  \begin{multline*}
    (d)_0=2\land\fn{DischargeLabel}(d)=0\land\fn{LastRule}(d)=1\land{}\\
    \fn{EndFmla}(d)={}\\
    \Gn{(}\concat\fn{EndFmla}((d)_1)\concat\Gn{\land}
    \concat\fn{EndFmla}((d)_2)\concat\Gn{)}.
  \end{multline*}

  మరొక సరళమైన ఉదాహరణ $\Intro\eq$ నియమం. దీనికి పూర్వాధారాలు లేవు కనుక,
  పరికల్పనల మాదిరిగానే $(d)_0=0$. దీనికి ఉపసంహరణ చీటీ కూడా లేదు;
  అంటే $n=0$. అయితే, ఏదో మూసిన పదం~$t$కి $!A$ అనేది $\eq[t][t]$ రూపంలో
  ఉండాలి. ఇక్కడ ఆదిమ పునరావృత్త నిర్వచనం
  \begin{multline*}
    (d)_0=0\land\fn{DischargeLabel}(d)=0\land{}\\
    \bexists{t<d}{(\fn{ClTerm}(t)\land
      \fn{EndFmla}(d)={}\\
      \Gn{{\eq}(}\concat t\concat\Gn{,}\concat t\concat\Gn{)})}.
  \end{multline*}

  మరింత క్లిష్టమైన ఉదాహరణగా, $\fn{FollowsBy}_{\Intro{\lif}}(d)$ ఎప్పుడు
  నిలుస్తుందో చూద్దాం. $\delta$ అంత్య-\tetoken{సూత్రం}{!!{formula}}
  $(!A\lif !B)$ రూపంలో ఉండి, $\delta_1$ అంత్య-\tetoken{సూత్రం}{!!{formula}}
  $!B$గా, $\delta$లో చీటీ~$n$ గల ప్రతి పరికల్పనా $!A$ రూపంలో ఉన్నప్పుడే
  అది నిలుస్తుంది. దీన్ని ఆదిమ పునరావృత్తంగా ఇలా వ్యక్తీకరించవచ్చు:
  \begin{multline*}
    (d)_0=1\land{}\\
    \bexists{a<d}{(\fn{Discharge}(a,(d)_1,\fn{DischargeLabel}(d))\land{}}\\
      \fn{EndFmla}(d)=(\Gn{(}\concat a\concat\Gn{\lif}
      \concat\fn{EndFmla}((d)_1)\concat\Gn{)}))
  \end{multline*}
  ($a$ను $!A$ గ్యోడెల్ సంఖ్యగా భావించండి.)

  ఇంకొక ఉదాహరణగా \Intro{\lexists}ను పరిగణించండి. ఇక్కడ,
  \tetoken{ఒక సూత్రం}{!!a{formula}}~$!A$, మూసిన పదం~$t$,
  \tetoken{ఒక చరం}{!!a{variable}}~$x$ ఉండి, $\Subst{!A}{t}{x}$ అనేది
  \tetoken{వ్యుత్పత్తి}{!!{derivation}}~$\delta_1$ అంత్య-
  \tetoken{సూత్రం}{!!{formula}}గా, $\lexists[x][!A]$ చివరి నిగమనపు
  నిష్కర్షగా ఉన్నప్పుడే $\delta$లోని చివరి నిగమనం సరైనది. కాబట్టి
  $\fn{FollowsBy}_{\Intro{\lexists}}(d)$ నిలవడానికి అవసరమైనదీ, సరిపడేదీ
  \begin{multline*}
    (d)_0=1\land\fn{DischargeLabel}(d)=0\land{}\\
    \bexists{a<d}{\bexists{x<d}{\bexists{t<d}{
          (\fn{ClTerm}(t)\land\fn{Var}(x)\land{}}}}\\
    \fn{Subst}(a,t,x)=\fn{EndFmla}((d)_1)\land
    \fn{EndFmla}(d)=(\Gn{\lexists}\concat x\concat a)).
  \end{multline*}

  తరువాత $\fn{Correct}(d)$ను ఇలా నిర్వచిస్తాం:
  \begin{multline*}
    \fn{Sent}(\fn{EndFmla}(d))\land{}\\
    [(\fn{LastRule}(d)=1\land
    \fn{FollowsBy}_{\Intro\land}(d))\lor\dots\lor{}\\
    (\fn{LastRule}(d)=16\land\fn{FollowsBy}_{\Elim\eq}(d))\lor{}\\
    \bexists{n<d}{\bexists{x<d}{(d=\tuple{0,x,n})}}].
  \end{multline*}
  \sourcecorrection{OLTEART-010}{ఆంగ్ల మూలంలోని $\fn{Correct}(d)$ సూత్రం మొత్తం వికల్పాన్ని కుండలీకరించలేదు; సాధారణ సంధి-వికల్ప ప్రాధాన్యతలో $\fn{Sent}$ పరీక్ష మొదటి నియమ సందర్భానికే వర్తిస్తుంది. వివరణ ఉద్దేశించినట్లుగా, ప్రతి నియమ సందర్భానికీ మరియు పరికల్పన సందర్భానికీ అంత్యసూత్రం వాక్యమని పరీక్షించేలా చతురస్ర కుండలీకరణలు చేర్చాం.}
  మొదటి పంక్తి $d$ అంత్య-\tetoken{సూత్రం}{!!{formula}} ఒక వాక్యమని
  నిర్ధారిస్తుంది. చివరి పంక్తి $d$ కేవలం పరికల్పన అయిన సందర్భాన్ని కలుపుతుంది.
\end{proof}

\begin{prob}
  \olref[inc][art][pnd]{prop:followsby}లోలాగే ఈ ధర్మాలను నిర్వచించండి:
  \begin{enumerate}
  \item $\fn{FollowsBy}_{\Elim{\lif}}(d)$,
  \item $\fn{FollowsBy}_{\Elim{\eq}}(d)$,
  \item $\fn{FollowsBy}_{\Elim{\lor}}(d)$,
  \item $\fn{FollowsBy}_{\Intro{\lforall}}(d)$.
  \end{enumerate}
  చివరిదానికి, గ్యోడెల్ సంఖ్య~$d$ గల
  \tetoken{వ్యుత్పత్తి}{!!{derivation}} యొక్క చివరి నిగమనం ఐగెన్ చరరాశి
  షరతును నెరవేరుస్తుందా, అంటే $\Intro{\lforall}$ నిగమనం యొక్క ఐగెన్
  చరరాశి~$a$ అనేది $d$ అంత్య-\tetoken{సూత్రం}{!!{formula}}లో గానీ,
  $d$లోని తెరచిన పరికల్పనలో గానీ కనిపించదా అని ఆదిమ పునరావృత్తంగా
  పరీక్షించగలమని కూడా చూపాలి. దీని కోసం
  \olref[inc][art][pnd]{prop:openassum}లోని ఆదిమ పునరావృత్త విధేయం
  $\fn{OpenAssum}$ను ఉపయోగించవచ్చు.
\end{prob}

\begin{prop}
  \ollabel{prop:deriv}
  $d$ సరైన \tetoken{వ్యుత్పత్తి}{!!{derivation}}~$\delta$ గ్యోడెల్ సంఖ్య
  అయితే నిలిచే $\fn{Deriv}(d)$ సంబంధం ఆదిమ పునరావృత్తం.
\end{prop}

\begin{proof}
  \tetoken{ఒక వ్యుత్పత్తి}{!!^a{derivation}}~$\delta$లోని ప్రతి నిగమనం
  ఒక నియమపు సరైన వర్తింపు అయితే, అంటే దాని ప్రతి
  ఉప-\tetoken{వ్యుత్పత్తులు}{!!{derivation}s} సరైన నిగమనంతో ముగిస్తే,
  $\delta$ సరైనది. కాబట్టి $\fn{Deriv}(d)$ నిలవడానికి అవసరమైనదీ, సరిపడేదీ
  \[
  \bforall{i<\len{\fn{SubtreeSeq}(d)}}{\fn{Correct}((\fn{SubtreeSeq}(d))_i)}
  \]
\end{proof}

\begin{prop}
  \ollabel{prop:openassum} గ్యోడెల్ సంఖ్య~$d$ గల
  \tetoken{వ్యుత్పత్తి}{!!{derivation}}~$\delta$లో
  \tetoken{ఉపసంహరించని}{!!a{undischarged}} పరికల్పన~$!A$ గ్యోడెల్
  సంఖ్య $z$ అయితే నిలిచే $\fn{OpenAssum}(z,d)$ సంబంధం ఆదిమ పునరావృత్తం.
\end{prop}

\begin{proof}
  ఒక పరికల్పన ఘటన చీటీ~$n$తో $\delta$లోని, ఉపసంహరణ చీటీ~$n$ గల నియమంతో
  ముగిసే ఉప-\tetoken{వ్యుత్పత్తి}{!!{derivation}}లో కనిపిస్తే అది
  \tetoken{ఉపసంహరించబడింది}{!!{discharged}}. కాబట్టి $\delta$లో దాని
  ఘటనలలో కనీసం ఒకటి \tetoken{ఉపసంహరించబడకపోతే}{!!{discharged}},
  $!A$ అనేది $\delta$లో \tetoken{ఉపసంహరించని}{!!a{undischarged}}
  పరికల్పన. ఇక్కడ జాగ్రత్త అవసరం: $\delta$లో $!A$కు
  \tetoken{ఉపసంహరించిన}{!!{discharged}},
  \tetoken{ఉపసంహరించని}{!!{undischarged}} ఘటనలు రెండూ ఉండవచ్చు.

  $\delta_0=\delta$, $\delta_k$ అనేది ఏదో~$n$కు పరికల్పన
  $\Discharge{!A}{n}$, ప్రతి $\delta_{i+1}$ అనేది $\delta_i$కి తక్షణ
  ఉప-\tetoken{వ్యుత్పత్తి}{!!{derivation}} అయ్యే $\delta_0$, \dots,
  $\delta_k$ క్రమాన్ని పరిగణించండి. ఏ $\delta_i$ కూడా ఉపసంహరణ చీటీ~$n$
  గల నిగమనంతో ముగియని అటువంటి క్రమం ఉంటే, $!A$ అనేది $\delta$లో
  \tetoken{ఉపసంహరించని}{!!a{undischarged}} పరికల్పన.

  ఆదిమ పునరావృత్త ప్రమేయం $\fn{SubtreeSeq}(d)$, $\delta$ యొక్క అన్ని
  ఉప-\tetoken{వ్యుత్పత్తుల}{!!{derivation}s} గ్యోడెల్ సంఖ్యల క్రమాన్ని
  ఇస్తుంది. $\delta$ ఉప-\tetoken{వ్యుత్పత్తుల}{!!{derivation}s} గ్యోడెల్
  సంఖ్యల ఏ క్రమమైనా దానికి ఉపక్రమం. ఉపక్రమంగా ఉండటం ఆదిమ పునరావృత్త
  సంబంధం: $\fn{Subseq}(s,s')$ నిలవడానికి అవసరమైనదీ, సరిపడేదీ
  $\bforall{i<\len{s}}{\lexists[j<\len{s'}][(s)_i=(s')_j]}$.
  తక్షణ ఉప-\tetoken{వ్యుత్పత్తి}{!!{derivation}}గా ఉండటం కూడా అలాంటిదే:
  $\fn{Subderiv}(d,d')$ నిలవడానికి అవసరమైనదీ, సరిపడేదీ
  $\bexists{j<(d')_0}{d=(d')_{j+1}}$.
  \sourcecorrection{OLTEART-011}{ఆంగ్ల మూలంలోని తక్షణ ఉపవ్యుత్పత్తి పరీక్ష $d=(d')_j$ను వాడుతుంది. కానీ $(d')_0$ అనేది పూర్వాధారాల సంఖ్య; ఉపవ్యుత్పత్తి సంకేతాలు స్థానాలు $1$ నుంచి $(d')_0$ వరకూ ఉంటాయి. అందువల్ల సూచికను $j+1$గా సరిచేశాం.}
  కాబట్టి $\fn{OpenAssum}(z,d)$ను ఇలా నిర్వచించవచ్చు:
  \begin{multline*}
    \bexists{s<\fn{SubtreeSeq}(d)}{(\fn{Subseq}(s,\fn{SubtreeSeq}(d))
      \land(s)_0=d\land{}}\\
    \bexists{n<d}{((s)_{\len{s}\tsub 1}=\tuple{0,z,n}\land{}}\\
      \bforall{i<(\len{s}\tsub 1)}{(\fn{Subderiv}((s)_{i+1},(s)_i)\land{}}\\
      \fn{DischargeLabel}((s)_i)\neq n))).
  \end{multline*}
\end{proof}

\begin{prop}
  \ollabel{prop:prf-prim-rec}
  $\Gamma$ అనేది \tetoken{వాక్యాల}{!!{sentence}s} ఆదిమ పునరావృత్త సమితి
  అనుకుందాం. ``$\Gamma$లోని \tetoken{ఉపసంహరించని}{!!{undischarged}}
  పరికల్పనల నుంచి $!A$ యొక్క \tetoken{ఒక వ్యుత్పత్తి}{!!a{derivation}}
  ~$\delta$ సంకేతసంఖ్య $x$, $!A$ గ్యోడెల్ సంఖ్య $y$'' అని వ్యక్తీకరించే
  $\Prf[\Gamma](x,y)$ సంబంధం ఆదిమ పునరావృత్తం.
\end{prop}

\begin{proof}
  ``$y\in\Gamma$''ను ఆదిమ పునరావృత్త విధేయం $R_\Gamma(y)$ ఇస్తుందని
  అనుకుందాం. $y$ ఒక వాక్యం~$!A$ గ్యోడెల్ సంఖ్యగా, $x$ అంత్య-
  \tetoken{సూత్రం}{!!{formula}}~$!A$ గల, \tetoken{ఉపసంహరించని}{!!{undischarged}}
  పరికల్పనలన్నీ $\Gamma$లో ఉండే సహజ నిగమన
  \tetoken{వ్యుత్పత్తి}{!!{derivation}} సంకేతసంఖ్యగా ఉన్నప్పుడే నిలిచే
  $\Prf[\Gamma](x,y)$ ఆదిమ పునరావృత్తమని చూపాలి.

  \olref{prop:deriv} ప్రకారం, $x$ సహజ నిగమనంలో సరైన
  \tetoken{వ్యుత్పత్తి}{!!{derivation}}~$\delta$ గ్యోడెల్ సంఖ్య అయితేనే
  నిలిచే $\fn{Deriv}(x)$ ధర్మం ఆదిమ పునరావృత్తం. కాబట్టి
  $\Prf[\Gamma](x,y)$ను ఇలా నిర్వచించవచ్చు:
  \begin{align*}
    \Prf[\Gamma](x,y)\defiff{}
    &\fn{Deriv}(x)\land\fn{EndFmla}(x)=y\land{}\\
    &\bforall{z<x}{(\fn{OpenAssum}(z,x)\lif R_\Gamma(z))}.
  \end{align*}
\end{proof}

\end{document}
